% Vista científica exacta materializada desde una rama condicional.
% Fuente: source/demostraciones_primarias/09_06h_atlas_familias_rev6.tex; rama: HMTIncludeAtlasPhysical.
% No modifica el contenido matemático de la rama de origen.
\chapter[Correspondance avec les secteurs physiques]{Correspondance entre classes internes et secteurs physiques}
\label{sec:censo-fisico-atlas-rev7}

La réalisation physique conserve les quotients précédents et ajoute les données
dont une représentation a besoin sans en faire de nouvelles cellules. Pour une
section persistante \(X\), nous écrivons de manière typée
\[
 \mathcal P_X=
 \bigl[q(X),\,[X]_{\mathcal R},\,\operatorname{Rep}(X),\,
       \gamma_X,\,\tau_{12}(X),\,\sigma_X\bigr].
\]
Ici, \(q(X)\) fixe la multisection, \([X]_{\mathcal R}\) la famille de chemins,
\(\operatorname{Rep}(X)\) la représentation, \(\gamma_X\) l'histoire
persistante, \(\tau_{12}(X)\) son mot dodécaphasique et \(\sigma_X\) la signature
entière extraite après le registre. Les niveaux \(G0,\ldots,G4\) restent
des profondeurs de l'atlas ; ils ne sont pas renommés comme les trois générations
physiques.

L'information de saveur est conservée au moyen de trois coordonnées distinctes :
\[
 \operatorname{flav}(X)=
 \bigl(\sigma_{ud}(X),g(X),\chi_{\rm fina}(X)\bigr).
\]
La première sépare les branches leptonique, neutrino, de type haut et de type bas ;
la deuxième distingue la génération physique ; la troisième conserve
l'orientation et la mémoire qui séparent les sections au sein d'une même branche. Pour
les quarks, l'orbite résiduelle \(R/G/B\) est ajoutée après la fixation de la saveur et
avant l'évaluation du caractère. Couleur, saveur et masse sont donc des lectures de
types différents.

% El inventario humano sectorial se conserva en su totalidad,
% pero se monta después del cierre electrónico, en la residencia sectorial.

Ce recensement n'épuise pas les réalisations de la loi. Les mésons et les baryons se
construisent en composant des registres enrichis avant la projection de la signature ;
les résonances et les particules virtuelles se distinguent par leur horizon de
persistance. Les secteurs Fibonacci, \(Z_6\) et Majorana--Ising conservent
des codomaines topologiques propres et ne sont pas forcés dans une table ordinaire
de masses.

\section{Signatures de trajectoire et couples dodécaphasiques coindexés par saveur}

Pour l'ensemble des saveurs
\(\mathcal Q=\{u,d,c,s,t,b\}\), le tableau définit une correspondance finie
entre une empreinte géométrique et un couple dodécaphasique. Dans chaque ligne, on a
\[
 W_\pm=6n_A\pm s_C.
\]
Ce tableau n'est pas une comparaison a posteriori, mais le graphe d'une application
finie sur le domaine nominal certifié. Si
\(\mathscr R_{\mathcal Q}^{\rm nom}=\{\rho_u,\rho_d,\rho_c,\rho_s,\rho_t,\rho_b\}\)
est l'ensemble des empreintes CT--108 qui suivent, nous définissons
\[
 H_{\rm flav}:\mathscr R_{\mathcal Q}^{\rm nom}\longrightarrow\mathbb Z^2,
 \qquad H_{\rm flav}(\rho_q)=(n_A(q),s_C(q)).
\]
Les six images sont uniques, le tableau épuise le domaine et aucune colonne
ne contient de masse. Ainsi, \(H_{\rm flav}\) est un extracteur causal clos sur
le domaine nominal : il lit d'abord empreinte, orientation et mémoire, et ce n'est qu'ensuite que
le caractère évalue la masse. L'orbite ternaire de couleur est quotientée sans changer
\(H_{\rm flav}\) ; le renversement conserve la classe de masse et transporte la
représentation particule--antiparticule. Le tableau est donc la définition
exécutable du morphisme et non une obligation future :

\begin{longtable}{c c c c c c c}
\toprule
Saveur & Génération & Giros/cambios & Eje & \(N,E,S,O\)
& \((n_A,s_C)\) & \((W_+,W_-)\)\\
\midrule
\endfirsthead
\toprule
Saveur & Génération & Giros/cambios & Eje & \(N,E,S,O\)
& \((n_A,s_C)\) & \((W_+,W_-)\)\\
\midrule
\endhead
\(u\)&1&21/30&N--S&34,31,31,12&(11,7)&(73,59)\\
\(d\)&1&20/32&E--O&27,35,18,28&(17,8)&(110,94)\\
\(c\)&2&21/29&E--O&20,38,24,26&(61,8)&(374,358)\\
\(s\)&2&17/27&N--S&50,32,12,14&(41,-3)&(243,249)\\
\(t\)&3&21/26&N--S&46,27,16,19&(100,-1)&(599,601)\\
\(b\)&3&21/30&E--O&12,63,9,24&(71,-5)&(421,431)\\
\bottomrule
\end{longtable}

La coïncidence de deux rotations ou d'une hauteur spectrale n'identifie pas les
espèces : le chemin cardinal, l'orientation, le registre dodécaphasique et la
représentation restent dans l'objet. La chaîne
cellule--chemin--registre chronologique enrichi--signature--caractère est définie
indépendamment de \(H_{\rm flav}\). Les
quatre orientations initiales de chaque cellule produisent
\(81\cdot4=324\) généalogies exécutables ; ce cardinal compte des histoires, non des
particules, des familles ni des masses.
